%% ma: L12-B14-0077
%% lop: 12
%% chuong: 5
%% bai: 14
%% dang: 12.14.23
%% loai: TN4
%% muc: TH
%% dap_an: "A"
%% nguon: "To_n_12_PT_M_t_Ph_ng.pdf (Đề kiểm tra đánh giá kiến thức, Chuyên đề Phương trình mặt phẳng, Nguyễn Anh Khôi - Lê Minh Khôi), Câu 32"
%% kien_thuc: "Bài 14 (vectơ pháp tuyến), Bài 8 (tọa độ vectơ)"
%% trang_thai: cho-duyet
%% ly_do: "Dạng 12.14.23 mới đề xuất, chờ giáo viên duyệt."
%% ghi_chu: "Tài liệu gốc không kèm đáp án; đáp án do Claude giải và kiểm bằng sympy."
\begin{ex}
Trong không gian $Oxyz$, cho hai điểm $M(1;2;-4)$ và $M'(5;4;2)$, biết $M'$ là hình chiếu vuông góc của $M$ lên mặt phẳng $(\alpha)$. Khi đó mặt phẳng $(\alpha)$ có một vectơ pháp tuyến là
\choice{\True $\vec n=(2;1;3)$}{$\vec n=(2;3;3)$}{$\vec n=(3;3;-1)$}{$\vec n=(2;-1;3)$}
\loigiai{$MM'\perp(\alpha)$ nên $(\alpha)$ nhận $\overrightarrow{MM'}=(4;2;6)=2(2;1;3)$ làm vectơ pháp tuyến. Chọn A.}
\end{ex}
